% Question bank for Lecture 17 (entity authentication). All numbers verified by script (scratchpad/w22/verify*.py).
\section{Entity Authentication}\label{q:entityauth}

\subsection*{Check your understanding}
\begin{enumerate}[leftmargin=1.6em,itemsep=3pt]
  \item \textbf{Define entity authentication. Who are the claimant and the verifier?} Entity authentication is a technique that lets one party prove the identity of another party. The entity whose identity must be proved is the claimant; the party that tries to prove the claimant's identity is the verifier. When Bob checks Alice, Alice is the claimant and Bob the verifier.
  \item \textbf{Give the two differences between data-origin (message) authentication and entity authentication.} First, message authentication need not happen in real time (Alice may be offline when Bob checks her e-mail), whereas entity authentication does (Alice must be online and take part before any message flows, as at an ATM). Second, message authentication covers one message and must be repeated for every message, whereas entity authentication covers the claimant for the whole session.
  \item \textbf{Name the three kinds of witness with two examples each. Is ``password plus a security question'' two-factor authentication?} Something known (password, PIN, private key), something possessed (smart card, passport, ATM card), something inherent (fingerprint, retina, voice, signature). Password plus security question is \emph{not} two-factor: both are something known, so one keylogger or one phishing page captures both; two factors must come from two different categories.
  \item \textbf{Why does a system store $h(P)$ instead of $P$? Which attack is still possible, and why do account lock-outs not stop it?} If the file of clear passwords is stolen, every account is lost at once; with hashes, Eve cannot invert $h$ to get $P$, and sending $h(P)$ does not help because the system hashes whatever it receives. The remaining attack is the dictionary attack: Eve hashes every likely password and looks for matches in the stolen file. It runs offline on Eve's own computer, so the system never sees her guesses; lock-outs and delays only limit \emph{online} guessing, and Eve goes online only once, with the password she has already found.
  \item \textbf{What does salting achieve, and what does it not achieve?} The salt, a random string stored in clear beside $h(P \| S)$, makes equal passwords hash differently, multiplies the size of any precomputed table by $2^b$ for a $b$-bit salt (4096 for UNIX's 12 bits), forces the users of one file to be attacked one at a time, and hides password reuse across systems. It does not slow an attack on one chosen user, because that user's salt is in the file; only strong passwords and slow (iterated) hashing help there.
  \item \textbf{What is the weakness of the second one-time-password approach (sequential update)?} Each password $P_{i+1}$ travels encrypted under $P_i$, so the chain can be followed forwards: if Eve guesses or steals $P_1$ (or any $P_i$), she decrypts the recorded traffic and learns all later passwords. The system must also keep the current $P_i$ in usable form, and a lost message desynchronises the two sides.
  \item \textbf{In Lamport's scheme, why can an eavesdropper not use the value she sees?} The values are used in the order $h^{n-1}(P_0), h^{n-2}(P_0), \dots$; after seeing $h^k(P_0)$ Eve needs $h^{k-1}(P_0)$, a preimage of what she holds, which a preimage-resistant hash does not allow. Replaying $h^k(P_0)$ fails because the system now stores it and checks $h(\text{received})$ against it.
  \item \textbf{What is the small-$n$ attack on Lamport's scheme and how is it prevented?} Eve impersonates the server and sends a counter $m$ much smaller than the real one; Alice answers $h^{m-1}(P_0)$, from which Eve can hash forwards to any later chain element the real server will ask for. Alice must refuse a counter that is not smaller than the last one she answered (and should authenticate the server).
  \item \textbf{Differentiate a replay attack and a reflection attack.} A replay attack resends a message recorded in an earlier run and is defeated by freshness (nonce, timestamp, sequence number). A reflection attack sends a party's own message or response back to it in a concurrent run; it is defeated by asymmetry: the claimant answers first, identities or different formats inside responses, or different keys per direction.
  \item \textbf{Compare a nonce and a timestamp as the challenge.} A nonce needs no synchronised clocks or shared state but costs an extra message from the verifier. A timestamp lets the claimant start, so one message suffices, but the clocks must be synchronised and the verifier must remember recent timestamps to stop replays inside the acceptance window.
  \item \textbf{In signature-based challenge--response, why does Alice add her own nonce $R_A$ and Bob's name?} Without $R_A$, Alice would sign any value Bob chooses, which a dishonest verifier could use to obtain her signature on a document of his choosing. The verifier's name prevents Eve from presenting a response signed for Bob to another verifier, Carol.
  \item \textbf{State the three properties of a zero-knowledge proof.} Completeness: an honest claimant who knows the secret is always accepted. Soundness: a claimant without the secret is accepted only with negligible probability, which repetition makes as small as desired. Zero-knowledge: the verifier learns nothing except that the claim is true, because he could have produced every transcript himself.
  \item \textbf{Why must Alice choose a fresh $r$ in every Fiat--Shamir round?} If she uses the same $r$ in two rounds and gets both $c = 0$ and $c = 1$, she reveals $y_0 = r$ and $y_1 = rs$, and then $s = y_1 y_0^{-1} \bmod n$.
  \item \textbf{Differentiate FAR and FRR. What is the equal error rate?} FAR is the fraction of impostor attempts that are wrongly accepted (a security failure); FRR is the fraction of genuine attempts that are wrongly rejected (an inconvenience). Raising the matching threshold lowers FAR and raises FRR; the threshold at which they are equal gives the equal error rate, a single figure for comparing systems.
  \item \textbf{Why can a biometric template not be hashed like a password, and what is the main disadvantage of biometrics?} Two readings of the same feature are never identical, so the system must compute a similarity score against a stored usable template rather than compare hashes for equality. The main disadvantage is that a stolen or copied biometric cannot be changed, unlike a password; privacy of the template database is a further concern.
\end{enumerate}

\subsection*{Practice problems}
\begin{enumerate}[leftmargin=1.6em,itemsep=4pt]
  \item In the Fiat--Shamir protocol, $n = 253$ ($= 11 \times 23$) and Alice's private key is $s = 21$. (a) Find her public key $v$. (b) She chooses $r = 30$. Compute the witness, and the response and Bob's check for $c = 1$ and for $c = 0$. \hfill (4 marks)\\
  \textbf{Solution.} (a) $\gcd(21, 253) = 1$. $v = 21^2 \bmod 253 = 441 - 253 = \mathbf{188}$.\\
  (b) $x = 30^2 \bmod 253 = 900 - 759 = \mathbf{141}$.\\
  $c = 1$: $y = r s = 30 \cdot 21 = 630 \bmod 253 = 630 - 506 = \mathbf{124}$. Bob: $y^2 = 15\,376 = 60 \cdot 253 + 196 \to 196$; $x v = 141 \cdot 188 = 26\,508 = 104 \cdot 253 + 196 \to 196$. Equal, round passed.\\
  $c = 0$: $y = r = \mathbf{30}$. Bob: $y^2 = 900 \bmod 253 = 141 = x = x v^0$. Round passed.\\
  \emph{Check}: $y^2 = r^2 s^{2c} = x v^c$ in both cases.

  \item With $n = 77$ and Alice's public key $v = 4$, Bob receives the witness $x = 15$, sends $c = 1$, and receives $y = 13$. Should he accept the round? What would have happened if he had sent $c = 0$? \hfill (2 marks)\\
  \textbf{Solution.} $y^2 = 169 \bmod 77 = 15$, while $x v^c = 15 \cdot 4 = 60$. Since $15 \ne 60$, Bob \textbf{rejects}. The claimant sent $y = r$ (indeed $13^2 \equiv 15 = x$), which is the correct answer to $c = 0$ only; had Bob sent $c = 0$, the check $y^2 = x$ would have passed. This is exactly a cheater who bet on $c = 0$.\\
  \emph{Check}: the honest answer for $c = 1$ is $y = rs = 13 \cdot 9 \bmod 77 = 40$, with $40^2 \bmod 77 = 60$.

  \item (a) How many Fiat--Shamir rounds make a cheater's success probability less than $10^{-6}$? Less than $10^{-9}$? (b) How many Feige--Fiat--Shamir rounds with $k = 4$ secrets give less than $10^{-6}$? (c) How many Guillou--Quisquater rounds with $e = 3$, and with $e = 65537$? \hfill (4 marks)\\
  \textbf{Solution.} (a) A cheater passes a round with probability $1/2$, so $k$ rounds with $(1/2)^k$. Need $2^k > 10^6$: $\log_2 10^6 = 19.93$, so $\mathbf{20}$ rounds ($2^{-20} = 9.54 \times 10^{-7}$). For $10^{-9}$: $\log_2 10^9 = 29.90$, so $\mathbf{30}$ rounds ($2^{-30} = 9.31 \times 10^{-10}$).\\
  (b) Per round $1/2^4 = 1/16$; need $16^t > 10^6$: $16^4 = 65\,536$, $16^5 = 1\,048\,576$, so $\mathbf{5}$ rounds ($16^{-5} = 9.54 \times 10^{-7}$, the same as 20 Fiat--Shamir rounds).\\
  (c) Per round $1/e$. For $e = 3$: $3^{12} = 531\,441 < 10^6 < 3^{13} = 1\,594\,323$, so $\mathbf{13}$ rounds. For $e = 65537$: $65537^2 = 4\,295\,098\,369 > 10^6$, so $\mathbf{2}$ rounds (probability $2.33 \times 10^{-10}$).\\
  \emph{Check}: rounds needed $= \lceil \log(1/\varepsilon) / \log(\text{choices per round}) \rceil$.

  \item In Feige--Fiat--Shamir, $n = 143$ and Alice's private keys are $[3, 5, 7]$, so her public keys are $[16, 103, 108]$. She chooses $r = 10$ and Bob sends $[0, 1, 1]$. Compute $x$, $y$ and Bob's check. \hfill (3 marks)\\
  \textbf{Solution.} Public keys: $v_i = (s_i^2)^{-1}$: $9 \cdot 16 = 144 \equiv 1$, $25 \cdot 103 = 2575 = 18 \cdot 143 + 1$, $49 \cdot 108 = 5292 = 37 \cdot 143 + 1$ \checkmark.\\
  $x = 10^2 \bmod 143 = \mathbf{100}$. With $c = [0, 1, 1]$: $y = r s_2 s_3 = 10 \cdot 5 \cdot 7 = 350 \bmod 143 = 350 - 286 = \mathbf{64}$.\\
  Bob computes $y^2 v_2 v_3$: $64^2 = 4096 = 28 \cdot 143 + 92 \to 92$; $92 \cdot 103 = 9476 = 66 \cdot 143 + 38 \to 38$; $38 \cdot 108 = 4104 = 28 \cdot 143 + 100 \to \mathbf{100} = x$. Accepted.\\
  \emph{Check}: $y^2 \prod v_i^{c_i} = x \prod (s_i^2 v_i)^{c_i} = x$.

  \item In Guillou--Quisquater, $n = 77$ ($= 7 \times 11$), $e = 7$, and Alice's secret is $s = 2$. (a) Verify that $e$ is valid and find Alice's public key $v$. (b) She picks $r = 3$ and Bob sends $c = 2$. Compute $x$, $y$ and Bob's check. \hfill (5 marks)\\
  \textbf{Solution.} (a) $\phi = 6 \times 10 = 60$; $\gcd(7, 60) = 1$, so $e = 7$ is valid. $s^e = 2^7 = 128 \bmod 77 = 51$. $v = 51^{-1} \bmod 77$ by extended Euclid on $(77, 51)$: quotients $1, 1, 1, 25$; remainders $26, 25, 1, 0$; $t$-values $-1, 2, -3$. So $v = -3 + 77 = \mathbf{74}$. Check: $51 \cdot 74 = 3774 = 49 \cdot 77 + 1$ \checkmark.\\
  (b) $x = 3^7 \bmod 77 = 2187 - 28 \cdot 77 = 2187 - 2156 = \mathbf{31}$. $y = r s^c = 3 \cdot 2^2 = \mathbf{12}$.\\
  Bob: $y^7$: $12^2 = 144 \to 67$; $12^4 = 67^2 = 4489 \to 23$; $12^7 = 12^4 \cdot 12^2 \cdot 12$: $23 \cdot 67 = 1541 \to 1$, $1 \cdot 12 = 12$. $v^2 = 74^2 = 5476 = 71 \cdot 77 + 9 \to 9$. $y^e v^c = 12 \cdot 9 = 108 \bmod 77 = \mathbf{31} = x$. Accepted.\\
  \emph{Check}: $y^e v^c = r^e (s^e v)^c = x \cdot 1^c$.

  \item Lamport's scheme uses the toy hash $h(x) = x^2 \bmod 1003$. Alice's original password is $P_0 = 25$ and $n = 4$. (a) Compute the hash chain and the entry stored by the system. (b) Show the first two logins and the stored entry after each. (c) Eve records the value sent in the first login. What can she compute from it, and why can she not log in? (d) How many logins are possible? \hfill (5 marks)\\
  \textbf{Solution.} (a) $h^1 = 25^2 = 625$; $h^2 = 625^2 = 390\,625 = 389 \cdot 1003 + 458 \to 458$; $h^3 = 458^2 = 209\,764 = 209 \cdot 1003 + 137 \to 137$; $h^4 = 137^2 = 18\,769 = 18 \cdot 1003 + 715 \to 715$. Stored entry: \textbf{(Alice, 4, 715)}.\\
  (b) Login 1: system sends 4; Alice sends $h^3(P_0) = 137$; system checks $137^2 \bmod 1003 = 715$ \checkmark; new entry \textbf{(Alice, 3, 137)}. Login 2: system sends 3; Alice sends $h^2(P_0) = 458$; $458^2 \bmod 1003 = 137$ \checkmark; new entry \textbf{(Alice, 2, 458)}.\\
  (c) From 137 Eve can only hash forwards: $h(137) = 715$, $h(715) = 698$, values that are already used or never asked. At login 2 she needs $h^2(P_0) = 458$, a square root of 137 modulo 1003, i.e.\ a preimage; replaying 137 gives $137^2 \bmod 1003 = 715 \ne 137$ and is rejected.\\
  (d) $n = 4$ logins (the last one sends $P_0 = 25$ itself); then the counter is 0 and the system must be re-initialised.

  \item A UNIX-style system uses a 12-bit salt. Eve's dictionary has $2^{20}$ words and she computes $2^{20}$ hashes per second. (a) How many entries does a precomputed table for all salts need, and how long does building it take? (b) The stolen file has $2^{10}$ users with distinct salts. How long does it take to test the dictionary against all of them, with and without salt? (c) How long for one targeted user? \hfill (4 marks)\\
  \textbf{Solution.} (a) $2^{20} \times 2^{12} = 2^{32} = \mathbf{4\,294\,967\,296}$ entries, $2^{12} = 4096$ times the unsalted table; building it takes $2^{32}/2^{20} = 2^{12} = 4096$ s $\approx \mathbf{1.14}$ \textbf{hours}.\\
  (b) Without salt: one pass of $2^{20}$ hashes, each looked up among the stored hashes: $\mathbf{1}$ \textbf{s}. With distinct salts: one pass per user, $2^{10} \times 2^{20} = 2^{30}$ hashes $= 1024$ s $\approx \mathbf{17}$ \textbf{minutes}.\\
  (c) Her salt is in the file, so one pass of $2^{20}$ hashes: $\mathbf{1}$ \textbf{s}, the same as without salt.\\
  \emph{Check}: salting multiplies the bulk attack by the number of distinct salts and precomputation by $2^{12}$, but not the single-user attack.

  \item Passwords are 8 characters drawn from the 62 letters and digits, and the stored hashes are unsalted. Eve can compute $10^9$ hashes per second. (a) How long does an exhaustive offline search take? (b) How much longer for 9 characters? (c) An ATM card is blocked after 3 wrong 4-digit PINs; what is a thief's chance of success? \hfill (3 marks)\\
  \textbf{Solution.} (a) $62^8 = 218\,340\,105\,584\,896 \approx 2.18 \times 10^{14}$ hashes; $2.18 \times 10^{14} / 10^9 = 2.18 \times 10^5$ s $\approx 60.7$ hours $\approx \mathbf{2.5}$ \textbf{days}.\\
  (b) One more character multiplies the work by 62: $\approx \mathbf{157}$ \textbf{days}.\\
  (c) Three distinct guesses out of $10^4$ PINs: $3/10^4 = \mathbf{0.03\%}$. The short PIN is safe only because guessing is online and limited; offline, $10^4$ hashes would take microseconds.

  \item A face-recognition system is tested with 2000 genuine attempts, of which 50 are rejected, and 8000 impostor attempts, of which 4 are accepted. (a) Compute FRR and FAR. (b) The system is used for identification against 1000 enrolled faces with this FAR per comparison. What is the probability that an impostor matches someone? (c) What happens to both rates if the threshold is raised? \hfill (3 marks)\\
  \textbf{Solution.} (a) $\mathrm{FRR} = 50/2000 = \mathbf{2.5\%}$; $\mathrm{FAR} = 4/8000 = \mathbf{0.05\%}$.\\
  (b) $1 - (1 - 0.0005)^{1000} = 1 - 0.6065 = \mathbf{0.394}$ (about $1 - e^{-0.5}$, since $1000 \times 0.0005 = 0.5$ false matches are expected).\\
  (c) A stricter threshold lowers FAR and raises FRR; the rates are equal at the equal error rate.

  \item Draw and explain challenge--response mutual authentication using a symmetric-key cipher. Show a reflection attack on the version in which the initiator challenges first, and state how to prevent it. \hfill (6 marks)\\
  \textbf{Solution (model answer outline).} Draw the four-message diagram of \S\ref{sec:entityauth}: (1) Alice; (2) $R_B$; (3) $E_{K_{AB}}(R_A \| R_B)$; (4) $E_{K_{AB}}(R_B \| R_A)$. Explain: message 3 proves Alice knows $K_{AB}$ (contains Bob's fresh nonce) and carries her challenge; message 4 proves Bob knows $K_{AB}$; fresh nonces stop replay. Then the vulnerable protocol: (1) Alice, $R_A$; (2) $E_K(R_A)$, $R_B$; (3) $E_K(R_B)$. Attack in two sessions: S1: Eve sends ``Alice'', $R_E$, gets $E_K(R_E)$, $R_B$; S2: Eve sends ``Alice'', $R_B$, gets $E_K(R_B)$; she returns $E_K(R_B)$ in S1 and is accepted. Prevention: claimant proves first; identities or different formats inside the ciphertext ($E_K(R_B, \text{Bob})$); different keys per direction.

  \item Explain the Fiat--Shamir zero-knowledge protocol with a diagram. Prove that an honest claimant always passes, and show that a dishonest claimant passes a round with probability $1/2$. \hfill (6 marks)\\
  \textbf{Solution (model answer outline).} Setup: trusted party publishes $n = pq$; Alice's secret $s$ (coprime to $n$), public $v = s^2 \bmod n$. Diagram: $x = r^2$ (witness) $\to$; $\leftarrow c \in \{0, 1\}$; $y = r s^c \to$; Bob checks $y^2 \equiv x v^c$. Proof: $y^2 = r^2 s^{2c} = x v^c$. Cheater: guessing $c = 1$, send $x = r^2 v^{-1}$ and $y = r$; guessing $c = 0$, send $x = r^2$ and $y = r$; each passes only if the guess is right. She cannot answer both for one $x$, since $y_0^2 = x$, $y_1^2 = xv$ give $(y_1/y_0)^2 = v$, a square root of $v$ (as hard as factoring $n$). Hence probability $1/2$ per round, $(1/2)^k$ after $k$ rounds (20 rounds: $9.54 \times 10^{-7}$). Mention the zero-knowledge property and the need for a fresh $r$.

  \item Explain biometric authentication: enrollment, verification and identification, the two classes of techniques with examples, and the two accuracy measures. \hfill (5 marks)\\
  \textbf{Solution (model answer outline).} Definition: measurement of physiological or behavioural features (something inherent). Components: capturing device, processor (feature extraction, template, matching), storage. Enrollment: each person's template recorded in the database before use. Verification: 1:1 match against the claimed identity's record (bank checks a signature on a cheque); identification: 1:N match against all records (employee-only building). Draw the enrollment/authentication flow with the threshold decision. Physiological: fingerprint, iris, retina, face, hands, voice, DNA; behavioural: signature, keystroke. Accuracy: FRR $=$ false rejections / genuine attempts; FAR $=$ false acceptances / impostor attempts; threshold trade-off and EER. Close with advantages (cannot be forgotten or lent) and disadvantages (cannot be changed if stolen, copying, privacy).
\end{enumerate}
